% Generated draft: requires mathematical and source review.
\documentclass[11pt]{article}
\usepackage[margin=25mm]{geometry}
\usepackage{fontspec}
\setmainfont{Latin Modern Roman}
\usepackage{amsmath,amssymb,mathtools,unicode-math}
\setmathfont{Latin Modern Math}
\usepackage{enumitem,needspace,xurl,hyperref}
\setlength{\parindent}{0pt}
\setlength{\parskip}{5pt}
\setlength{\emergencystretch}{3em}
\newcommand{\N}{\mathbb{N}}
\newcommand{\Z}{\mathbb{Z}}
\newcommand{\Q}{\mathbb{Q}}
\newcommand{\R}{\mathbb{R}}
\newcommand{\C}{\mathbb{C}}
\newcommand{\F}{\mathbb{F}}
\newcommand{\A}{\mathbb{A}}
\newcommand{\PP}{\mathbb{P}}
\newcommand{\E}{\mathbb{E}}
\newcommand{\eps}{\varepsilon}
\newcommand{\dd}{\,\mathrm{d}}
\begin{document}
\section*{UnsolvedMath ID: 800052}
\textbf{Source problem code:} catalogue entry
\subsection{Definitions and mathematical statement}
For an integer \(k\geq 0\), let \(M\) be a closed smooth oriented manifold of real dimension \(4k\), where closed means compact without boundary. Its signature \(\tau(M)\) is the number of positive eigenvalues minus the number of negative eigenvalues of the nondegenerate symmetric bilinear form
\[
H^{2k}(M;\R)\times H^{2k}(M;\R)\longrightarrow\R,
\qquad (a,b)\longmapsto\langle a\smile b,[M]\rangle,
\]
where \([M]\) is the fundamental class determined by the orientation. Write \(p_i(TM)\in H^{4i}(M;\Z)\) for the Pontryagin classes of the tangent bundle, and also for their images in rational cohomology. The Hirzebruch \(L\)-class is the multiplicative characteristic class characterized, in formal Pontryagin roots, by
\[
p(E)=\prod_j(1+x_j^2),\qquad
L(E)=\prod_j\frac{x_j}{\tanh x_j}.
\]
Let \(L_k(p_1,\ldots,p_k)\) denote its component of cohomological degree \(4k\); in particular, \(L_0=1\) and \(L_1=p_1/3\).

\textbf{Requested proof.} Give an elementary proof that, for every such \(M\),
\[
\tau(M)=\bigl\langle L_k\bigl(p_1(TM),\ldots,p_k(TM)\bigr),[M]\bigr\rangle.
\]
This is the Hirzebruch signature theorem, called the index theorem in the historical problem. The supplied account attributes its proof to Hirzebruch at the beginning of June 1953; the requested contribution is an elementary proof of that established identity. The source does not specify which methods qualify as elementary.

\textbf{Additional requested explanation.} Explain the geometrical meaning of the functional equation valid for \(\tau\). The supplied record does not state this functional equation or specify its variables, hypotheses, or domain, so this part cannot be formulated precisely from the supplied evidence.

\textbf{Dimensional qualification.} The source describes \(M\) as any closed smooth oriented manifold without explicitly giving its dimension. The displayed numerical signature formula is formulated here in dimension \(4k\), using the degree-\(4k\) component of the \(L\)-class. The record does not specify conventions or an intended extension in other dimensions.
\subsection{Short English statement}
Find an elementary proof that the signature of every closed smooth oriented \(4k\)-manifold equals its top-degree Hirzebruch \(L\)-class evaluated on its fundamental class. Also explain the geometry behind the functional equation mentioned in the source. That equation and the intended meaning of elementary are unspecified. The supplied registry labels the problem open; this is not an independently verified assessment, and the signature identity itself is presented as an established theorem.
\subsection{Sources}
\begin{itemize}
\item Local UnsolvedMath registry, supplied complete record: permanent ID \(800052\), exact source problem code \texttt{catalogue entry}, title ``Hirzebruch's 1954 Problem List --- Problem 1.'' The supplied status is \texttt{open}; no independent status verification is asserted.
\item Kotschick, \emph{Updates on Hirzebruch's 1954 Problem List} (2013), Problem 1. Supplied source URL: \url{https://arxiv.org/abs/1305.4623}. The provenance records access on 2026-10-16 and extraction by \texttt{arxiv\_latex}; these are supplied provenance statements, not a claim that the source was visited or independently checked for this draft.
\item Attribution: the registry lists Kotschick as proposer and 2013 as the proposed year, while identifying the item with Hirzebruch's 1954 problem list. The supplied historical account attributes the original signature theorem to Hirzebruch at the beginning of June 1953 and distinguishes its historical name, index theorem, from the later Fredholm-index terminology associated with the Atiyah--Singer index theorem.
\end{itemize}
\end{document}
